\chapter{Actor-Critic Methods}

REINFORCE (9.8) weights each $\nabla_\theta\log\pi_\theta(A_t\mid S_t)$ by $\gamma^tG_t$, the full return of Chapter 4. That weight has two problems:

\begin{enumerate}
\item \textbf{No comparison.} $G_t$ is an absolute number, not a comparison: it says how good the continuation after $A_t$ was, not whether $A_t$ was better or worse than the other actions available in $S_t$.
\item \textbf{Full-return noise.} $G_t$ sums every reward from $t$ to the end of the episode, so it varies a lot from one episode to the next, and so does each update. The update must also wait until the episode ends before every $G_t$ is known.
\end{enumerate}
Sections 10.1--10.2 are \textbf{Monte Carlo}: they keep the return $G_t$ and fix problem 1. Section 10.3 is \textbf{TD}: it replaces $G_t$ by a one-step backup, as Chapter 5 did for values, and fixes problem 2.

Problem 1 comes first. For example, suppose that in $S_t$ the current policy picks $a_1$ or $a_2$ with probability $\tfrac12$ each, and that $a_1$ leads to a return of $100$ and $a_2$ to a return of $200$. When $a_1$ is sampled, (9.8) moves $\theta$ by $\alpha\,100\gamma^t\,\nabla_\theta\log\pi_\theta(a_1\mid S_t)$. The gradient points in the direction that increases $\log\pi_\theta(a_1\mid S_t)$, and its weight $100\gamma^t$ is positive, so the step raises the probability of $a_1$ above $\tfrac12$, and that of $a_2$ falls below $\tfrac12$, although $a_1$ is the worse choice. When $a_2$ is sampled, the same kind of step, with weight $200\gamma^t$, raises the probability of $a_2$. Each update favors whichever action was sampled, but an update for $a_2$ is twice as large, so the policy moves toward $a_2$ only on average, while every sample of $a_1$ pushes the wrong way. The fix is to subtract from each return the average return from $S_t$ under the current policy, $\tfrac12\cdot100+\tfrac12\cdot200=150$. Then $a_1$ gets weight $-50\gamma^t$ and $a_2$ gets weight $+50\gamma^t$, so every sample of $a_1$ lowers its probability and every sample of $a_2$ raises it. A quantity subtracted from the return in this way is called a \textbf{baseline}.

The average return from $S_t$ under the current policy is, by definition, the state value $V^{\pi_\theta}(S_t)=\mathbb{E}[G_t\mid S_t]$ (B.7); in the example, $V^{\pi_\theta}(S_t)=150$. It is unknown, so it is learned as an estimate $\hat{V}(s;\mathbf{w})$, as in Chapter 7. Subtracting $\hat{V}(S_t;\mathbf{w})$ from $G_t$ gives \textbf{REINFORCE with baseline} (Section 10.1), with weight $G_t-\hat{V}(S_t;\mathbf{w})$. Section 10.2 names that comparison the \textbf{advantage}. Problem 2 remains because the weight still uses $G_t$. Section 10.3 replaces $G_t-\hat{V}(S_t;\mathbf{w})$ by the TD error $R_{t+1}+\gamma\hat{V}(S_{t+1};\mathbf{w})-\hat{V}(S_t;\mathbf{w})$, the same substitution Chapter 5 made for values. Used this way, $\hat{V}$ is called the \textbf{critic}, because it judges each action by comparing what happened after one step with what it expected from $S_t$. The policy $\pi_\theta$, which chooses the actions and is updated by policy gradient as in Chapter 9, is the \textbf{actor}. A method in which the critic's judgment is the weight in the actor's update is an \textbf{actor--critic} method (\citetxt{barto1983}{Barto, Sutton, and Anderson, 1983}). REINFORCE with baseline also learns $\hat{V}$, but uses it only as a baseline, so it is not usually called actor--critic.

\section{Baseline}

Like REINFORCE and the Monte Carlo methods of Chapter 4, this section uses the full return $G_t$. It addresses problem 1: replace that absolute return by a comparison against a state-dependent baseline. The example above subtracted the average return $150$ from every return in $S_t$. For that fix to be valid, the update must still estimate $\nabla_\theta J$ after the subtraction. Start from (9.7) and subtract from each return a \textbf{baseline} $b(s)$, a function of the state alone. Then

\[
\begin{aligned}
&\mathbb{E}_{\tau\sim\pi_\theta}\Biggl[
\sum_{t=0}^{T-1}
\gamma^{t}\bigl(G_t-b(S_t)\bigr)\,\nabla_\theta\log\pi_\theta(A_t\mid S_t)
\Biggr]
\\
&\qquad=
\mathbb{E}_{\tau\sim\pi_\theta}\Biggl[
\sum_{t=0}^{T-1}
\gamma^{t}G_t\,\nabla_\theta\log\pi_\theta(A_t\mid S_t)
\Biggr]
-
\sum_{t=0}^{T-1}\gamma^{t}\,
\mathbb{E}\bigl[
b(S_t)\,\nabla_\theta\log\pi_\theta(A_t\mid S_t)
\bigr]
\quad\text{[A.4]}
\\
&\qquad=
\nabla_\theta J(\theta)
-
\sum_{t=0}^{T-1}\gamma^{t}\cdot 0
\quad\text{[(9.7) and (E.3)]}
\\
&\qquad=
\nabla_\theta J(\theta).
\end{aligned}
\]

So subtracting the baseline leaves the gradient unchanged:

\begin{equation}
\nabla_\theta J(\theta)
=
\mathbb{E}_{\tau\sim\pi_\theta}\Biggl[
\sum_{t=0}^{T-1}
\gamma^{t}\bigl(G_t-b(S_t)\bigr)\,\nabla_\theta\log\pi_\theta(A_t\mid S_t)
\Biggr].
\tag{10.1}
\end{equation}

The $G_t$ term stays because $G_t$ depends on the action $A_t$. For the baseline term, apply the tower property (A.8) by conditioning on $S_t$: \[
\begin{aligned}
\mathbb{E}\bigl[b(S_t)\,\nabla_\theta\log\pi_\theta(A_t\mid S_t)\bigr]
&=
\mathbb{E}\Bigl[\mathbb{E}\bigl[b(S_t)\,\nabla_\theta\log\pi_\theta(A_t\mid S_t)\bigm|S_t\bigr]\Bigr]
\\
&=
\mathbb{E}\Bigl[b(S_t)\,\mathbb{E}\bigl[\nabla_\theta\log\pi_\theta(A_t\mid S_t)\bigm|S_t\bigr]\Bigr]
\\
&=
\mathbb{E}\bigl[b(S_t)\cdot 0\bigr]
=
0.
\end{aligned}
\] Given $S_t$, the baseline is a fixed number, so it factors out of the inner average over $A_t$; that average is zero by \eqref{eq:E.1}. The same display is \eqref{eq:E.3} in \appref{E}.

The choice of $b$ therefore changes how noisy the estimate is, not its mean. A good choice is the one in the example, the average return from $S_t$, which is $b(s)=V^{\pi_\theta}(s)=\mathbb{E}[G_t\mid S_t=s]$. With it, the weight $G_t-V^{\pi_\theta}(S_t)$ is positive when $A_t$ did better than the policy's average in $S_t$ and negative when it did worse. This baseline is learned as $\hat{V}(s;\mathbf{w})$, a second approximator with its own weights $\mathbf{w}$, as in Chapter 7, and (10.1) holds for every $\mathbf{w}$, even before $\hat{V}$ is accurate.

REINFORCE with baseline therefore learns two things from the same episodes: the value $\hat{V}(\cdot;\mathbf{w})$ and the policy $\pi_\theta$. Each has its own parameters and its own objective, so each episode gives two updates. The weights $\mathbf{w}$ take a gradient-descent step that brings $\hat{V}$ closer to the observed returns. The parameters $\theta$ take a gradient-ascent step on $J$, as in REINFORCE, with $\hat{V}$ as the baseline. The two updates are linked only through that baseline: the policy update reads the current value of $\hat{V}(S_t;\mathbf{w})$ as a number and does not change $\mathbf{w}$.

Learning $\hat{V}$ is the prediction problem of Chapter 4: $V^{\pi_\theta}(s)=\mathbb{E}[G_t\mid S_t=s]$, so each sampled return $G_t$ is a target for $\hat{V}(S_t;\mathbf{w})$. Gradient descent on the squared error $\tfrac12\bigl(G_t-\hat{V}(S_t;\mathbf{w})\bigr)^2$, summed over the episode, gives

\begin{equation}
\mathbf{w}
\leftarrow
\mathbf{w}
+
\alpha_{\mathbf{w}}
\sum_{t=0}^{T-1}
\bigl(G_t-\hat{V}(S_t;\mathbf{w})\bigr)\,
\nabla_{\mathbf{w}}\hat{V}(S_t;\mathbf{w}).
\tag{10.2}
\end{equation}

The policy update is the \textbf{single-episode} estimate of (10.1), as (9.8) was for (9.7):

\begin{equation}
\theta
\leftarrow
\theta
+
\alpha_\theta
\sum_{t=0}^{T-1}
\gamma^{t}\bigl(G_t-\hat{V}(S_t;\mathbf{w})\bigr)\,
\nabla_\theta\log\pi_\theta(A_t\mid S_t).
\tag{10.3}
\end{equation}

Comparing (10.2) and (10.3) with gradient descent gives the losses. In $L(\theta)$, both $G_t$ and $\hat{V}(S_t;\mathbf{w})$ are fixed numbers:

\begin{equation}
L(\mathbf{w})
=
\frac12\sum_{t=0}^{T-1}
\bigl(G_t-\hat{V}(S_t;\mathbf{w})\bigr)^2,
\qquad
L(\theta)
=
-\sum_{t=0}^{T-1}
\gamma^{t}\bigl(G_t-\hat{V}(S_t;\mathbf{w})\bigr)\,
\log\pi_\theta(A_t\mid S_t).
\tag{10.4}
\end{equation}

Both updates use the same differences $G_t-\hat{V}(S_t;\mathbf{w})$, with $\hat{V}$ evaluated at the $\mathbf{w}$ from before this episode's value step. That keeps the baseline a function of $S_t$ alone. If $\mathbf{w}$ were first adjusted to match this episode's $G_t$ and then $\hat{V}$ were recomputed for (10.3), the baseline would depend on the episode's actions $A_t$, and $\mathbb{E}\bigl[\hat{V}(S_t;\mathbf{w})\,\nabla_\theta\log\pi_\theta(A_t\mid S_t)\bigr]$ need not be zero. The order of the two parameter steps does not matter once those differences are fixed.

\begin{implnote}[Algorithm (REINFORCE with baseline)]
Input: differentiable policy $\pi_\theta(a\mid s)$; differentiable state-value function $\hat{V}(s;\mathbf{w})$; step sizes $\alpha_\theta$, $\alpha_{\mathbf{w}}$; discount $\gamma$; number of episodes $M$.

Output: trained policy $\pi_\theta$ and $\hat{V}$.

\begin{enumerate}
\item Initialize $\theta$ and $\mathbf{w}$ (e.g. randomly).
\item Loop for episode $=1,2,\ldots,M$:
\begin{itemize}
\item Generate an episode $\tau=(S_0,A_0,R_1,\ldots,S_T)$ by sampling $A_t\sim\pi_\theta(\cdot\mid S_t)$ until termination.
\item For each $t=0,1,\ldots,T-1$:
\begin{itemize}
\item $G_t\leftarrow\sum_{k=t}^{T-1}\gamma^{k-t}R_{k+1}$.
\end{itemize}
\item Evaluate $\hat{V}(S_t;\mathbf{w})$ for each $t$ (before updating $\mathbf{w}$).
\item $\mathbf{w}\leftarrow\mathbf{w}+\alpha_{\mathbf{w}}\sum_{t=0}^{T-1}\bigl(G_t-\hat{V}(S_t;\mathbf{w})\bigr)\nabla_{\mathbf{w}}\hat{V}(S_t;\mathbf{w})$.
\item $\theta\leftarrow\theta+\alpha_\theta\sum_{t=0}^{T-1}\gamma^{t}\bigl(G_t-\hat{V}(S_t;\mathbf{w})\bigr)\nabla_\theta\log\pi_\theta(A_t\mid S_t)$.
\end{itemize}
\item Return $\pi_\theta$ and $\hat{V}$.
\end{enumerate}

\end{implnote}

\lstinputlisting{../code/ch10/reinforce_baseline.py}


\section{Advantage function}

Section 10.1 chose the Monte Carlo baseline $b(s)=V^{\pi_\theta}(s)$, so the policy update weights each action by $G_t-V^{\pi_\theta}(S_t)$: how that action's return compares with the average return from the state. In the opening example, that weight was $-50\gamma^t$ for $a_1$ and $+50\gamma^t$ for $a_2$. Naming that difference makes the comparison itself the object of study: the policy gradient can be written with it as the weight, and Section 10.3 can replace this Monte Carlo estimate by a TD one (problem 2).

The weight $G_t-V^{\pi_\theta}(S_t)$ is random. Given $S_t=s$ and $A_t=a$, expand its conditional expectation by linearity (A.5): \[
\begin{aligned}
&\mathbb{E}\bigl[G_t-V^{\pi_\theta}(S_t)\bigm|S_t=s,A_t=a\bigr]
\\
&\qquad=
\mathbb{E}[G_t\mid S_t=s,A_t=a]
-
\mathbb{E}\bigl[V^{\pi_\theta}(S_t)\bigm|S_t=s,A_t=a\bigr].
\end{aligned}
\] The first term is $Q^{\pi_\theta}(s,a)$ by definition (B.13). In the second, $S_t=s$ is given, so $V^{\pi_\theta}(S_t)$ is the fixed number $V^{\pi_\theta}(s)$, and the expectation of a constant is that constant. Hence \[
\mathbb{E}\bigl[G_t-V^{\pi_\theta}(S_t)\bigm|S_t=s,A_t=a\bigr]
=
Q^{\pi_\theta}(s,a)-V^{\pi_\theta}(s).
\] The \textbf{advantage} of action $a$ in state $s$ under policy $\pi$ is that difference: how much is gained by taking $a$ compared with the average return $V^{\pi}(s)$ from the state $s$.

\begin{equation}
A^{\pi}(s,a)
=
Q^{\pi}(s,a)-V^{\pi}(s).
\tag{10.5}
\end{equation}

The expectation of the advantage, when the action is drawn from $\pi$ in $s$, is zero. Take the expectation of both sides of (10.5) and apply linearity (A.4): \[
\begin{aligned}
\mathbb{E}_{A\sim\pi(\cdot\mid s)}\bigl[A^{\pi}(s,A)\bigr]
&=
\mathbb{E}_{A\sim\pi(\cdot\mid s)}\bigl[Q^{\pi}(s,A)-V^{\pi}(s)\bigr]
\\
&=
\mathbb{E}_{A\sim\pi(\cdot\mid s)}\bigl[Q^{\pi}(s,A)\bigr]
-
\mathbb{E}_{A\sim\pi(\cdot\mid s)}\bigl[V^{\pi}(s)\bigr].
\end{aligned}
\] The first term equals $V^{\pi}(s)$ by (B.19). The second equals $V^{\pi}(s)$ as well: given $s$, $V^{\pi}(s)$ is a constant, so its expectation is itself. The two terms cancel, and the expectation of the advantage is zero. Under $\pi$, better-than-average and worse-than-average actions cancel. In the example, $A^{\pi_\theta}(S_t,a_1)=100-150=-50$ and $A^{\pi_\theta}(S_t,a_2)=200-150=+50$, and $\tfrac12(-50)+\tfrac12(+50)=0$.

The same comparison appears in the $Q$ form of the policy gradient. Equation (9.7a) weights $\nabla_\theta\log\pi_\theta(A_t\mid S_t)$ by $Q^{\pi_\theta}(S_t,A_t)$. Subtracting the baseline $V^{\pi_\theta}(S_t)$, as in Section 10.1, replaces that weight by $Q^{\pi_\theta}(S_t,A_t)-V^{\pi_\theta}(S_t)=A^{\pi_\theta}(S_t,A_t)$:

\begin{equation}
\nabla_\theta J(\theta)
=
\mathbb{E}_{\tau\sim\pi_\theta}\Biggl[
\sum_{t=0}^{T-1}
\gamma^{t}A^{\pi_\theta}(S_t,A_t)\,\nabla_\theta\log\pi_\theta(A_t\mid S_t)
\Biggr].
\tag{10.6}
\end{equation}

REINFORCE with baseline estimates the advantage by the Monte Carlo error $G_t-\hat{V}(S_t;\mathbf{w})$, where $\hat{V}(\cdot;\mathbf{w})$ is the learned approximation of $V^{\pi_\theta}$ from Section 10.1. That estimate is unbiased when $\hat{V}=V^{\pi_\theta}$, because $\mathbb{E}[G_t\mid S_t,A_t]=Q^{\pi_\theta}(S_t,A_t)$. It still uses the full return $G_t$, so problem 2 remains.

\section{One-step actor-critic}

Sections 10.1--10.2 fixed problem 1, but the advantage is still estimated with the full return $G_t$ (problem 2). Start from the advantage itself (10.5):

\[
A^{\pi_\theta}(s,a)
=
Q^{\pi_\theta}(s,a)-V^{\pi_\theta}(s).
\]

Substitute (B.20) for $Q^{\pi_\theta}$:

\[
\begin{aligned}
A^{\pi_\theta}(s,a)
&=
\mathbb{E}_\pi\bigl[R_{t+1}+\gamma V^{\pi_\theta}(S_{t+1})\mid S_t=s,A_t=a\bigr]
-
V^{\pi_\theta}(s)
\\
&=
\mathbb{E}_\pi\bigl[R_{t+1}+\gamma V^{\pi_\theta}(S_{t+1})-V^{\pi_\theta}(s)\mid S_t=s,A_t=a\bigr].
\end{aligned}
\]

So $A^{\pi_\theta}(s,a)$ is already that conditional expectation: a one-step TD error of the true $V^{\pi_\theta}$. Sample one transition and replace the unknown $V^{\pi_\theta}$ by the learned $\hat{V}(\cdot;\mathbf{w})$:

\begin{equation}
\delta_t
=
R_{t+1}+\gamma\hat{V}(S_{t+1};\mathbf{w})-\hat{V}(S_t;\mathbf{w}),
\tag{10.7}
\end{equation}

with $\hat{V}(S_{t+1};\mathbf{w})=0$ when $S_{t+1}$ is terminal. Write $U_t=R_{t+1}+\gamma\hat{V}(S_{t+1};\mathbf{w})$, so $\delta_t=U_t-\hat{V}(S_t;\mathbf{w})$---the same TD error as (5.3). If $\hat{V}=V^{\pi_\theta}$, then $\delta_t$ is an unbiased sample of $A^{\pi_\theta}(S_t,A_t)$ by the display above; if not, $\delta_t$ is biased. In exchange it depends on one reward and one next-state value (problem 2).

Section 10.2 wrote the policy gradient with the advantage as the weight (10.6). Estimate that weight by $\delta_t$:

\begin{equation}
\nabla_\theta J(\theta)
\approx
\mathbb{E}_{\tau\sim\pi_\theta}\Biggl[
\sum_{t=0}^{T-1}
\gamma^{t}\delta_t\,\nabla_\theta\log\pi_\theta(A_t\mid S_t)
\Biggr].
\tag{10.7a}
\end{equation}

Here $\hat{V}(S_{t+1};\mathbf{w})$ stands in for the rest of the return, so $\delta_t$ judges $A_t$: positive when the step beat what $\hat{V}$ expected from $S_t$, negative when it fell short. That judgment makes $\hat{V}$ the \textbf{critic} and $\pi_\theta$ the \textbf{actor}---an \textbf{actor--critic} method. In Section 10.1 the same $\hat{V}$ only subtracted from the observed $G_t$; every later reward stayed in the weight, so $\hat{V}$ was a baseline, not a critic.

After each transition, update both with the same $\delta_t$ (evaluated before the $\mathbf{w}$ step). The critic uses semi-gradient TD(0) on $U_t$ (Chapter 7); the actor uses the one-step form of (10.3):

\begin{equation}
\mathbf{w}
\leftarrow
\mathbf{w}
+
\alpha_{\mathbf{w}}\,
\delta_t\,
\nabla_{\mathbf{w}}\hat{V}(S_t;\mathbf{w}),
\tag{10.8}
\end{equation}

\begin{equation}
\theta
\leftarrow
\theta
+
\alpha_\theta\,
\gamma^{t}\delta_t\,
\nabla_\theta\log\pi_\theta(A_t\mid S_t).
\tag{10.9}
\end{equation}

With $U_t$ and $\delta_t$ treated as fixed numbers, the losses are

\begin{equation}
L(\mathbf{w})
=
\frac12\bigl(U_t-\hat{V}(S_t;\mathbf{w})\bigr)^2,
\qquad
L(\theta)
=
-\gamma^{t}\delta_t\,\log\pi_\theta(A_t\mid S_t).
\tag{10.10}
\end{equation}

Compute $\delta_t$ with the current $\mathbf{w}$ before either parameter step. The algorithm below updates $\mathbf{w}$ first, then $\theta$, both with that same $\delta_t$. As in Section 10.1, once $\delta_t$ is fixed the order of the two steps does not matter.

\begin{implnote}[Algorithm (one-step actor--critic)]
Input: differentiable policy $\pi_\theta(a\mid s)$; differentiable state-value function $\hat{V}(s;\mathbf{w})$; step sizes $\alpha_\theta$, $\alpha_{\mathbf{w}}$; discount $\gamma$; number of episodes $M$.

Output: trained policy $\pi_\theta$ and $\hat{V}$.

\begin{enumerate}
\item Initialize $\theta$ and $\mathbf{w}$ (e.g. randomly).
\item Loop for episode $=1,2,\ldots,M$:
\begin{itemize}
\item Initialize $S$ (the start state $S_0$); $t\leftarrow 0$.
\item Loop until $S$ is terminal:
\begin{itemize}
\item Sample $A\sim\pi_\theta(\cdot\mid S)$; take $A$, observe $R$, $S'$.
\item $\delta\leftarrow R+\gamma\hat{V}(S';\mathbf{w})-\hat{V}(S;\mathbf{w})$, with $\hat{V}(S';\mathbf{w})=0$ if $S'$ is terminal.
\item $\mathbf{w}\leftarrow\mathbf{w}+\alpha_{\mathbf{w}}\,\delta\,\nabla_{\mathbf{w}}\hat{V}(S;\mathbf{w})$.
\item $\theta\leftarrow\theta+\alpha_\theta\,\gamma^{t}\delta\,\nabla_\theta\log\pi_\theta(A\mid S)$.
\item $S\leftarrow S'$; $t\leftarrow t+1$.
\end{itemize}
\end{itemize}
\item Return $\pi_\theta$ and $\hat{V}$.
\end{enumerate}

\end{implnote}

\lstinputlisting{../code/ch10/actor_critic.py}

